\section{Introduction}

\label{sec:intro}
%% =====================================================================

The quantum threat to elliptic curve cryptography is not a question of
whether but when.  Shor's algorithm~\cite{shor1994} solves the elliptic
curve discrete logarithm problem in polynomial time on a quantum
computer.  Once a CRQC exists, every ECDSA private key is recoverable
from its public key, and every historical signature becomes forgeable.

For cryptocurrencies with short-lived state (e.g., daily trading
balances), the quantum threat can be mitigated by migration at the
time CRQCs become imminent.  For digital securities, this strategy
fails for three reasons:

\begin{enumerate}[nosep]
  \item \textbf{Long asset lifetimes.}  Bonds, real estate tokens, and
        structured products have maturities of 10--30 years.  A security
        issued in 2026 with a 30-year maturity must remain unforgeable
        until 2056.
  \item \textbf{Harvest-now, decrypt-later.}  An adversary can record
        signed transactions today and forge them once a CRQC is
        available.  For securities, this means retroactive forgery of
        ownership transfers.
  \item \textbf{Regulatory record-keeping.}  SEC Rule 17a-4 requires
        broker-dealers to retain records for 6 years; FINRA Rule 4511
        requires 6 years.  If signatures become forgeable during the
        retention period, the audit trail is compromised.
\end{enumerate}

The Lux network addresses this by deploying post-quantum cryptography
as native EVM precompiles at genesis.  This is not a future roadmap
item; it is a day-one capability.

%% =====================================================================
